\section*{第 \ref{ch:b1:electrostatics-gauss} 章——静电学：场与高斯定理}

\begin{solution}{exo:b1:electrostatics-gauss:1}
$F_e = 8.99 \times 10^9 \times (1.6 \times 10^{-19})^2/(2 \times 10^{-15})^2 = \qty{58}{N}$;
$F_g = 6.67 \times 10^{-11} \times (1.67 \times 10^{-27})^2/(4 \times 10^{-30}) =
\qty{4.6e-35}{N}$；比值 $1.2 \times 10^{36}$。原子核之所以能保持稳定，
强相互作用的作用范围在飞米量级，足以胜过这一斥力。
\end{solution}

\begin{solution}{exo:b1:electrostatics-gauss:2}
$E = 8.99 \times 10^9 \times 10^{-6}/1 = \qty{9.0}{kV/m}$。向下的场
是负电荷的场：$Q = -4\pi\varepsilon_0R^2E = -1.11 \times 10^{-10} \times
(6.37 \times 10^6)^2 \times 100 = \qty{-4.5e5}{C}$.
\end{solution}

\begin{solution}{exo:b1:electrostatics-gauss:3}
在轴上，$|x| > a$ 时：$E_x = \dfrac{q}{4\pi\varepsilon_0}\Bigl[\dfrac{1}{(x - a)^2} +
\dfrac{1}{(x + a)^2}\Bigr]$，方向向外。在 $y$ 轴上，$x$ 分量
相互抵消：$E_y = \dfrac{2qy}{4\pi\varepsilon_0(y^2 + a^2)^{3/2}}$；$\dd E_y/\dd y \propto
(y^2 + a^2) - 3y^2 = 0$，解得 $y = a/\sqrt2$。原点处的场
为零；从每个电荷射向另一电荷的\omterm{def:b1:electrostatics-gauss:field}{场线}，沿
$\pm y$ 方向偏折——形成鞍点。
\end{solution}

\begin{solution}{exo:b1:electrostatics-gauss:4}
两电荷之间两场均指向 $-q$，无零点。位于
$+2q$ 左侧，$2q/x^2 > q/(x + a)^2$ 恒成立。在 $-q$ 外侧，距其 $d$
处：$2q/(a + d)^2 = q/d^2$，$a + d = \sqrt2\,d$，$d = a/(\sqrt2 - 1) = 2.4a$。
从 $+2q$ 发出的两条\omterm{def:b1:electrostatics-gauss:field}{场线}，对应一条终止于 $-q$ 的\omterm{def:b1:electrostatics-gauss:field}{场线}；其余一半
延伸至无穷远。远处，这一对电荷看起来像电荷 $+q$。
\end{solution}

\begin{solution}{exo:b1:electrostatics-gauss:5}
半径为 $a$、宽度为 $\dd a$ 的圆环：$\dd q = 2\pi\sigma a\,\dd a$，$\dd E_z =
\dfrac{\sigma z}{2\varepsilon_0}\dfrac{a\,\dd a}{(a^2 + z^2)^{3/2}}$；积分得：
$\dfrac{\sigma z}{2\varepsilon_0}\Bigl[-\dfrac{1}{\sqrt{a^2 + z^2}}\Bigr]_0^R = \dfrac{\sigma}{2\varepsilon_0}
\Bigl(1 - \dfrac{z}{\sqrt{z^2 + R^2}}\Bigr)$。当 $z \ll R$ 时，$\sigma/2\varepsilon_0$，
即无限大平面的场。当 $z \gg R$ 时：$z/\sqrt{z^2 + R^2} \approx 1 - R^2/2z^2$，故 $E_z
\approx \sigma R^2/4\varepsilon_0z^2 = Q/4\pi\varepsilon_0z^2$，即点电荷的场。
\end{solution}

\begin{solution}{exo:b1:electrostatics-gauss:6}
$Q = 79e = \qty{1.26e-17}{C}$; $E(R) = 8.99 \times 10^9 \times 1.26 \times 10^{-17}/
(7 \times 10^{-15})^2 = \qty{2.3e21}{V/m}$；在 $R/2$ 处减半，为 \qty{1.2e21}{V/m}；在
\qty{1}{pm} 处：$1.14 \times 10^{-7}/10^{-24} = \qty{1.1e17}{V/m}$——$10^{15}$ 倍
空气的击穿场强，此距离正是电子所在之处。
\end{solution}

\begin{solution}{exo:b1:electrostatics-gauss:7}
半径为 $r$、长为 $h$ 的圆柱面：$2\pi rh\,E = \lambda h/\varepsilon_0$，$E =
\lambda/2\pi\varepsilon_0r$。电晕出现的条件为 $\lambda = 2\pi\varepsilon_0 rE = 2\pi \times 8.85 \times
10^{-12} \times 0.01 \times 3 \times 10^6 = \qty{1.7}{\micro C/m}$；在 \qty{5}{m} 处，
$E = 3 \times 10^6 \times 0.01/5 = \qty{6}{kV/m}$.
\end{solution}

\begin{solution}{exo:b1:electrostatics-gauss:8}
$E = \sigma/\varepsilon_0 = 10^{-6}/8.85 \times 10^{-12} = \qty{1.1e5}{V/m}$。一块极板
对自身不产生合力（其自身场的贡献
两两抵消）；另一极板在其所在处的场为 $\sigma/2\varepsilon_0$，故
$F = Q\sigma/2\varepsilon_0 = 10^{-6} \times 10^{-6}/(2 \times 8.85 \times 10^{-12}) =
\qty{5.6e-2}{N}$，为吸引力；压强为 $\sigma^2/2\varepsilon_0 = \qty{0.056}{Pa}$。
\end{solution}

\begin{solution}{exo:b1:electrostatics-gauss:9}
当 $r < a$ 时：内导体的电荷位于其表面，
内部无电荷，$E = 0$。当 $a < r < b$ 时：$E = \lambda/2\pi\varepsilon_0r$。当 $r > b$ 时：
包围的电荷为 $\lambda - \lambda = 0$，$E = 0$——电缆对外
不产生静\omterm{def:b1:electrostatics-gauss:field}{电场}。在 $r = a$ 处：$2 \times 8.99 \times 10^9 \times 10^{-8}/(5 \times 10^{-4}) =
\qty{3.6e5}{V/m}$.
\end{solution}

\begin{solution}{exo:b1:electrostatics-gauss:10}
$\rho = 3Q/4\pi(R_2^3 - R_1^3)$。当 $r < R_1$ 时：$E = 0$。当 $R_1 < r < R_2$ 时：$4\pi r^2E =
\rho\,\tfrac43\pi(r^3 - R_1^3)/\varepsilon_0$，故 $E = \dfrac{Q}{4\pi\varepsilon_0r^2}\,
\dfrac{r^3 - R_1^3}{R_2^3 - R_1^3}$。当 $r > R_2$ 时：$Q/4\pi\varepsilon_0r^2$。在 $R_1$ 处两侧
均给出 $0$；在 $R_2$ 处两侧也给出 $Q/4\pi\varepsilon_0R_2^2$。在壳层内 $E
\propto r - R_1^3/r^2$ 增大，故最大值位于 $R_2$ 处，此后
场按 $1/r^2$ 衰减。在空腔内场为零：置于其中的电荷
不受力——无论它处于何处。
\end{solution}

\begin{solution}{exo:b1:electrostatics-gauss:11}
(a) 地核 $\tfrac43\pi(3.48 \times 10^6)^3 \times 11000 = \qty{1.94e24}{kg}$；地幔
$\tfrac43\pi[(6.37 \times 10^6)^3 - (3.48 \times 10^6)^3] \times 4500 = \qty{4.08e24}{kg}$;
总质量 \qty{6.0e24}{kg}，误差在 $1\%$ 以内。(b) 表面 $GM/R^2 = \qty{9.9}{m/s^2}$；
地核边界处 $6.67 \times 10^{-11} \times 1.94 \times 10^{24}/(3.48 \times 10^6)^2 =
\qty{10.7}{m/s^2}$。(c) 取 $g = GM(r)/r^2$，$\dd M = 4\pi r^2\rho_m\,\dd r$：
$\dd g/\dd r = 4\pi G\rho_m - 2GM(r)/r^3 = 4\pi G[\rho_m - \tfrac23\bar\rho(r)]$,
其中 $\bar\rho(r) = M(r)/\tfrac43\pi r^3$ 为半径 $r$ 以内的平均密度。
$g$ 随深度增加（$\dd g/\dd r < 0$）当且仅当 $\rho_m < \tfrac23\bar\rho$。紧贴
地表之下，$\bar\rho = 5500$，$\tfrac23\bar\rho = 3700 < 4500$：$g$
先略有下降；在地核边界处 $\bar\rho = 11000$，$7300 >
4500$：此处升至 $10.7$。
\end{solution}

\begin{solution}{exo:b1:electrostatics-gauss:12}
(a) 球内，$E = er/4\pi\varepsilon_0R^3$，方向向外；作用于电子的 $\vect F =
-\dfrac{e^2}{4\pi\varepsilon_0R^3}\,\vect r$：这是弹簧式的力。(b) $\omega = \sqrt{e^2/4\pi\varepsilon_0mR^3}
= \sqrt{2.31 \times 10^{-28}/(9.11 \times 10^{-31} \times 10^{-30})} = \qty{1.6e16}{rad/s}$,
$f = \qty{2.5e15}{Hz}$，$\lambda = c/f = \qty{120}{nm}$。(c) 惊人地
接近 \qty{122}{nm} 谱线：原子的大小与强度都正确。
但该模型只给出一个频率，而非谱线系，且原子核
并非 \qty{0.1}{nm} 的小球，而是线度为 $10^{-15}$~m 的点——按经典理论，
绕核运动的电子会在纳秒量级内把能量辐射殆尽。
修正办法见\cref{ch:b1:quantum-introduction}。
\end{solution}

\begin{solution}{pb:b1:electrostatics-gauss:1}
\textbf{1.} $\vect F = -e\vect E$：向上。地表处向下的场
是负电荷产生的场：地球带负电。

\textbf{2.} 半径为 $R_E$ 的球面，$E_n = -100$：$Q = 4\pi\varepsilon_0R_E^2E_n =
1.11 \times 10^{-10} \times 4.06 \times 10^{13} \times (-100) = \qty{-4.5e5}{C}$.

\textbf{3.} $\sigma = \varepsilon_0E_n = \qty{-8.9e-10}{C/m^2}$; $8.9 \times 10^{-10}/1.6 \times
10^{-19} = 5.5 \times 10^9$ 个电子每平方米，$5.5 \times 10^5$ 个
每平方厘米。

\textbf{4.} $F = 100e \times 100 = \qty{1.6e-15}{N}$；重力 $10^{-9} \times 9.8 =
\qty{9.8e-9}{N}$：约为前者的六百万倍。

\textbf{5.} 外法线方向：顶面 $E_n = -5$，底面 $E_n = +100$，侧面
为零：$\Phi = \qty{95}{V.m}$，故 $Q = 95\varepsilon_0 = \qty{+8.4e-10}{C}$ 每平方
米——正的空间电荷，几乎抵消了地面的
$-8.9 \times 10^{-10}$.

\textbf{6.} $\rho = 8.4 \times 10^{-10}/10^4 = \qty{8.4e-14}{C/m^3}$: $5 \times 10^5$
\omterm{def:b1:electrostatics-gauss:charge}{元电荷}；与每种符号各 $10^9$ 个离子相比，
失衡度为 $5 \times 10^{-4}$。

\textbf{7.} $J = \gamma E_0 = \qty{2e-12}{A/m^2}$; $I = J \times 5.1 \times 10^{14} =
\qty{1.0}{kA}$；$Q/I = 4.5 \times 10^5/1000 = \qty{450}{s}$。地球
将在八分钟内变为电中性：必有某种过程在持续为它补充电荷。

\textbf{8.} 轴线位于包含它的每个平面之中，这些平面都是\omterm{prop:b1:electrostatics-gauss:symmetry}{对称面}：
$\vect E$ 沿轴向。每个 $\dd q$ 在其自身方向上给出 $\dd q/4\pi\varepsilon_0(a^2 + z^2)$
沿其自身方向，其中轴向部分所占比例为
$z/\sqrt{a^2 + z^2}$: $E_z = qz/4\pi\varepsilon_0(a^2 + z^2)^{3/2}$.

\textbf{9.} $\dd q = 2\pi\sigma a\,\dd a$: $E_z = \dfrac{\sigma z}{2\varepsilon_0}\displaystyle\int_0^R
\dfrac{a\,\dd a}{(a^2 + z^2)^{3/2}} = \dfrac{\sigma}{2\varepsilon_0}\Bigl(1 - \dfrac{z}{\sqrt{z^2
+ R^2}}\Bigr)$.

\textbf{10.} $z \ll R$：$\sigma/2\varepsilon_0$，即无限大平面（由\omterm{thm:b1:electrostatics-gauss:gauss}{高斯定理}得到）；
$z \gg R$：$\sigma R^2/4\varepsilon_0z^2 = Q/4\pi\varepsilon_0z^2$，即点电荷。

\textbf{11.} $\sigma = -40/(\pi \times 9 \times 10^6) = \qty{-1.4e-6}{C/m^2}$; $\sigma/2\varepsilon_0 =
\qty{8.0e4}{V/m}$；在 $z = \qty{3}{km}$ 处，$1 - 3/\sqrt{18} = 0.29$：$E =
\qty{2.3e4}{V/m}$，方向\emph{向上}，指向负电荷。

\textbf{12.} 上方圆盘，$z = \qty{9}{km}$：$8.0 \times 10^4 \times (1 - 9/\sqrt{90}) =
\qty{4.1e3}{V/m}$，方向向下。合场向上，为 $\qty{1.9e4}{V/m}$，
计入地面感应电荷后约为 $4 \times 10^4$，与实测的 \qty{10}{kV/m} \omterm{def:b1:units-dimensions:oom}{数量级}相当
（实际电荷分布更分散，且部分被屏蔽）；
方向相反，且约为晴天\omterm{def:b1:electrostatics-gauss:field}{电场}的一百倍。

\textbf{13.} 整朵云：净电荷为零，\omterm{def:b1:electrostatics-gauss:flux}{通量}为零。仅考虑下方圆盘：
$-Q/\varepsilon_0 = -40/8.85 \times 10^{-12} = \qty{-4.5e12}{V.m}$.

\textbf{14.} 每个圆盘相距 \qty{3}{km}，各产生 $2.3 \times 10^4$；
上方（$+$）圆盘向下推，下方（$-$）圆盘向下拉：\qty{4.7e4}{V/m}
方向向下——比击穿场强低三十倍。

\textbf{15.} $Q/4\pi\varepsilon_0R^2 = 8.99 \times 10^9 \times 40/(2.5 \times 10^5) =
\qty{1.4e6}{V/m}$：只有压缩在半公里范围内的电荷才能达到
击穿，而实测的云内\omterm{def:b1:electrostatics-gauss:field}{电场}约小十倍。局部
增强发生在雨滴与冰晶处，或宇宙射线引发的失控电子，
都是候选机制——闪电如何开始，仍有争议。

\textbf{16.} $I = 5/(5 \times 10^{-5}) = \qty{1e5}{A}$.

\textbf{17.} $W = QEd = 5 \times 10^5 \times 3000 = \qty{1.5e9}{J}$ (\qty{420}{kW h});
$P = 1.5 \times 10^9/5 \times 10^{-5} = \qty{3e13}{W}$——相当于全人类
总功率的一点五倍，持续仅五十微秒。

\textbf{18.} 体积 $\pi(0.02)^2 \times 3000 = \qty{3.8}{m^3}$，质量 \qty{4.5}{kg}；
$\Delta T = 1.5 \times 10^9/(4.5 \times 720) = \qty{4.6e5}{K}$，若无任何能量逸出。
通道向外爆炸（激波即雷声），并辐射出光
和无线电波，使空气电离、解离；把最终
通道加热到 $3 \times 10^4$~K 只占其中很小一部分。

\textbf{19.} 每千米 \qty{2.9}{s}。声音同时从一条 \qty{3}{km} 长的通道各点发出，
传播距离相差数千米，因此声音
在数秒内陆续到达——这就是隆隆声。

\textbf{20.} $50 \times 5 = \qty{250}{A}$ 的电流被导入大地；
加上较弱的雷暴电流（云下尖端放电、带电雨）后，总电流
达到第~I~部分的千安量级：雷暴是
全球电路的发电机，其回路为晴天空气。

\textbf{21.} $\Phi_{\text{外}}(\vect g) = -4\pi GM_{\text{内}}$：对点质量
$g = -GM/r^2$，球面面积为 $4\pi r^2$，\omterm{def:b1:electrostatics-gauss:flux}{通量}为 $-4\pi GM$，与
$r$ 无关；其余同库仑情形。

\textbf{22.} $g(r) = g_0r/R_E$；在 $R_E/2$ 处：\qty{4.9}{m/s^2}。

\textbf{23.} $M_{\text{地核}} = \qty{1.94e24}{kg}$，$g = GM_{\text{地核}}/R_c^2 =
\qty{10.7}{m/s^2}$——大于地表处的值：上方的地幔
比其下方物质的平均值更轻。

\textbf{24.} $m\ddot x = -mg_0x/R_E$: $\omega = \sqrt{g_0/R_E} = \qty{1.24e-3}{rad/s}$,
$T = 2\pi/\omega = \qty{5060}{s} = \qty{84}{min}$, $v_{\max} = \omega R_E =
\qty{7.9}{km/s}$——即贴地卫星的\omterm{thm:b1:central-forces:circular}{轨道速率}，其
周期也相同。

\textbf{25.} 它取代了对地球上每个电荷的求和（第~I~部分），
对空气中离子的求和（第~I~部分），以及对每层岩石的求和（第~IV~部分）——
而在第~II~部分中，它确定了圆盘求和的上下限。需保留：\qty{-5e5}{C} 与
\qty{1}{kA}，对应晴天的地球；\qty{40}{C} 与 \qty{e4}{V/m}，对应雷暴之下；
\qty{e5}{A} 与 \qty{e9}{J}，每次雷击；\qty{10.7}{m/s^2}，位于地核
边界处。
\end{solution}
